\section{成立条件と残された課題}
\label{sec:conditions-future-work}

本論文で証明した理想HHLの正しさには，$A$が可逆なエルミート行列であり，
固有値を正確に符号化でき，制御回転と逆位相推定を正確に実行できるという仮定がある．
有限精度では，固有値$\lambda$の近似誤差が逆数$1/\lambda$の誤差へ拡大し，
小さい固有値ほど影響が大きい．したがって条件数$\kappa$は，
成功確率だけでなく必要な推定精度にも関係する．

計算量を評価するためには，アルゴリズム本体の回路だけでなく，
入力状態$\ket b$の準備，制御された$e^{\mathrm{i}At}$の実行，
固有値から回転角を得る可逆計算，振幅増幅の反復を含める必要がある．
これらを効率よく実行できる行列と入力に対象を限ることが，
HHLの計算上の利点を論じる前提となる．

また，出力は解ベクトルの全成分ではなく量子状態$\ket x$である．
全成分を古典的に読み出すことを目的とすると，多数の測定が必要になる．
したがって，特定の観測量や射影確率を求める場合，または解状態を
後続の量子処理へ渡す場合のように，量子状態のまま利用できる問題設定が重要となる．

今後の課題として，具体的な行列アクセス模型の下での時間発展の実装，
状態準備の計算量，位相推定と制御回転を合わせた誤差評価，
および解状態から必要な量を読み出す測定方法を詳しく検討する必要がある．

